\documentclass[10pt,a4paper]{amsart}
\usepackage[margin=19mm]{geometry}
\usepackage{amsmath,amssymb,mathtools}
\usepackage[scr=rsfs,cal=euler]{mathalfa}
\usepackage{xfrac}
\usepackage[unicode,hidelinks,bookmarks=false]{hyperref}
\newcommand{\ie}{i.e.\ }
\newcommand{\cf}{cf.\ }
\newcommand{\resp}{respectively\ }
\newcommand{\Subsection}{\subsection}
\providecommand{\nobreakdash}{\nobreak\hbox{-}}
\setlength{\emergencystretch}{2em}
\multlinegap0pt
\usepackage{bm}
\usepackage{bbm}
\usepackage{dsfont}
\usepackage{xspace}
\usepackage{cancel}
\usepackage{mathtools}
\usepackage{amsthm}
\makeatletter
\@ifundefined{thm}{\newtheorem{thm}{Thm}}{}
\@ifundefined{remark}{\newtheorem{remark}[thm]{Remark}}{}
\makeatother
\def\Ac{{\mathcal A}}
\def\Alg{{\mathsf{Alg}}}
\def\A{{\mathbf A}}
\def\I{{\mathbf I}}
\def\Lc{{\mathcal L}}
\def\b{{\mathbf b}}
\def\x{{\mathbf x}}
\def\zero{{\mathbf 0}}
\title[Fine-grained Analysis and Faster Algorithms for Iteratively ]{Fine-grained Analysis and Faster Algorithms for Iteratively Solving Linear Systems}
\author{Micha{\l} Derezi\'nski, Daniel LeJeune, Deanna Needell, Elizaveta Rebrova}
\date{Source: arXiv:2405.05818v2 (CC BY 4.0)}
\begin{document}
\maketitle

\noindent\emph{Source and licence.} Fine-grained Analysis and Faster Algorithms for Iteratively Solving Linear Systems. Version arXiv:2405.05818v2 (CC BY 4.0), distributed under the Creative Commons Attribution 4.0 International licence, as recorded in the arXiv licence metadata for this exact version and shown on the abstract page https://arxiv.org/abs/2405.05818v2. Full rights notice: licenses/arxiv-2405.05818-CC-BY-4.0.txt.\par\smallskip

\noindent\textit{Editorial scope.} Counted unit: the Remark whose source label is \texttt{None} in the article (source0.tex lines 1327--1329, source label \texttt{None}), delivered together with the article's own complete original proof of it (source0.tex lines 1330--1372, one \texttt{proof} environment of 43 lines). No mathematics is written, repaired or re-derived: statement and proof are the article's own text line for line. Required context retained uncounted: none. Every definition, display and label of those lines is retained unchanged. Omitted from this package and not redistributed: 1--1326, 1373--1533. Nothing inside a retained excerpt is omitted or altered: every retained body is delivered exactly as the article's lines are written, so no omission receipt of any kind accompanies this package. Citations: the retained excerpts carry no citation command at all, so no bibliography entry of the article is transcribed and no reference is redistributed. Reference adaptation: none was needed and none was made; every reference inside the delivered unit resolves inside it, so no unresolved reference or citation is produced. Printed slips kept literal: no printed word, symbol, hypothesis or number of the counted statement or of its proof was altered. Layout and class: the article's own class is \texttt{article} with packages including geometry, inputenc, amsmath, amssymb, amsthm, xcolor, dsfont, hyperref, algorithm, algpseudocode, soul, tikz, pgfplots, lastpage; the wrapper is the lane's pinned \texttt{amsart} editorial wrapper at 10pt on A4, which restates the theorem-like environments the retained text needs and adds the article's own macro definitions that the retained text needs (\texttt{\textbackslash A}, \texttt{\textbackslash Ac}, \texttt{\textbackslash Alg}, \texttt{\textbackslash I}, \texttt{\textbackslash Lc}, \texttt{\textbackslash b}, \texttt{\textbackslash x}, \texttt{\textbackslash zero}), with extra wrapper packages (bm, bbm, dsfont, xspace, cancel, mathtools). The delivered numbering is the unit's own. Assets: no retained span contains a figure environment, a graphics-inclusion command, a PDF-page inclusion, a table or a TikZ picture; no image file is copied into this package, the package declares no asset dependency and no image file is redistributed with it. Licence: version arXiv:2405.05818v2 of this article is distributed under the Creative Commons Attribution 4.0 International licence, as recorded in the arXiv licence metadata for this exact version and shown on the abstract page https://arxiv.org/abs/2405.05818v2. A full rights notice is delivered at licenses/arxiv-2405.05818-CC-BY-4.0.txt. Characters: every retained excerpt is pure ASCII, so the delivered engine, XeLaTeX with its default Latin Modern text font, reports no missing character.
  \begin{remark}
    An analogous reduction can be obtained along the same lines if we replace $\Lc$ with the family of all square linear systems (dropping the positive definiteness).
  \end{remark}

\begin{proof}
First, let us use $\Lc_\ell(n,\kappa_\ell)$ to denote the above defined family of $n\times n$ linear systems from $\Lc$ with bounded spectral tail condition number. We break the proof down into two cases, depending on which of the two terms dominates the value of the max in the lower~bound.
\smallskip

\noindent    
\textbf{Case 1:}  $T_{\Ac,\Lc}(\ell,\infty,\epsilon,\delta) \geq T_{\Ac,\Lc}(n-\ell,\kappa_\ell,\epsilon,\delta)$. We  proceed via a proof by contradiction. Suppose that some $\Alg\in\Ac$ has query complexity less than $T_{\Ac,\Lc}(\ell,\infty,\epsilon,\delta)$ for solving $\Lc_\ell(n,\kappa_\ell)$. We are going to show how to use $\Alg$ to solve \emph{all} $\ell\times \ell$ linear systems in $\Lc$, thereby obtaining a contradiction. Suppose that $(\A,\b,\x_0)$ is an $\ell\times \ell$ linear system in $\Lc$, and define:
\begin{align*}
(\bar\A,\bar\b,\bar\x_0)=\bigg(\begin{bmatrix}
    \A&\zero\\
    \zero& \sigma_{\ell}(\A)\I_{n-\ell}
\end{bmatrix}, 
\begin{bmatrix}
    \b\\ 
    \zero
\end{bmatrix},
\begin{bmatrix}
    \x_0\\ 
    \zero
\end{bmatrix}
\bigg).
\end{align*}
Note that $(\bar\A,\bar\b,\bar\x_0)\in\Lc_\ell(n,\kappa_\ell)$, since $\bar\A$ is positive definite and $\frac{\sigma_{\ell}(\bar\A)}{\sigma_n(\bar\A)}=1\leq\kappa_\ell$. Moreover, if $\Alg$ runs on this system and returns $\bar\x$ such that $\|\bar\x-\bar\x^*\|_{\bar\A}^2\leq\epsilon\|\bar\A\|\|\bar\x_0-\bar\x^*\|^2$, where $\bar\x^*=\bar\A^{-1}\bar\b$, then the vector $\tilde\x$ consisting of the first $\ell$ coordinates of $\bar\x$ satisfies $\|\tilde\x-\x^*\|_{\A}^2\leq\epsilon\|\A\|\|\x_0-\x^*\|^2$, where $\x^*=\A^{-1}\b$. This gives us the contradiction.
\smallskip

\noindent
\textbf{Case 2:} $T_{\Ac,\Lc}(\ell,\infty,\epsilon,\delta) > T_{\Ac,\Lc}(n-\ell,\kappa_\ell,\epsilon,\delta)$. This case follows similarly. Suppose that some $\Alg\in\Ac$ has query complexity less than $T_{\Ac,\Lc}(n-\ell,\kappa_\ell,\epsilon,\delta)$ for solving $\Lc_\ell(n,\kappa_\ell)$. Now, consider an $(n-\ell)\times  (n-\ell)$ linear system task $(\A,\b,\x_0)\in\Lc$ with condition number $\kappa(\A)\leq \kappa_\ell$, and define:
\begin{align*}
(\bar\A,\bar\b,\bar\x_0)=\bigg(\begin{bmatrix}
    \sigma_1(\A)\I_{\ell}&\zero\\
    \zero& \A
\end{bmatrix}, 
\begin{bmatrix}
    \zero\\ 
    \b
\end{bmatrix},
\begin{bmatrix}
    \zero\\ 
    \x_0
\end{bmatrix}
\bigg).
\end{align*}
Note that, as in Case 1, we have $(\bar\A,\bar\b,\bar\x_0)\in\Lc_\ell(n,\kappa_\ell)$, since $\bar\A$ is positive definite and $\frac{\sigma_{\ell}(\bar\A)}{\sigma_n(\bar\A)} = \frac{\sigma_1(\A)}{\sigma_{n-\ell}(\A)}= \kappa(\A)\leq \kappa_\ell$. Now, solving this system to $\epsilon$ accuracy allows us to recover an $\epsilon$ approximate solution for $(\A,\b,\x_0)$, which is a contradiction with the definition of $T_{\Ac,\Lc}(n-\ell,\kappa_\ell,\epsilon,\delta)$. This concludes the proof of the lemma.
\end{proof}

\end{document}
